 \Chapter{CONCLUSION}\label{sec:Conclusion}
Cette section donnera un survol des constats dressés au fil du mémoire. La section commence avec une synthèse des travaux, puis enchaine avec les limitaions et des améliorations futures qui pourraient être apportées pour améliorer l'état des connaissances sur

%%
%%  SYNTHESE DES TRAVAUX / SUMMARY OF WORKS
%%
\section{Synthèse des travaux}
Le mémoire ci-dessus a documenté un outil web permettant d'enregistrer l'évolution de la réglementation d'urbanisme en matière de stationnement au fil du temps dans l'actuelle ville de Québec. Ces règlements ont ensuite été utilisés pour inférer la capacité de stationnement hors-rue, en utilisant les données du rôle foncier disponible en données ouvertes et les données cadastrales obtenues sur le portail universitaire. \par 
L'outil en question a été développé utilise une architecture web qui permet de créer un interface assez spécialisé pour gérer et manipuler l'ensemble de données assez complexe qui est nécessaire pour mettre en place la méthodologie. Des ébauches de méthodes de validation statistique et d'équivalence ont été mises en places mais les résultats n'ont pas encore été mis en place. 

%%
%%  LIMITATIONS
%%
\section{Limitations de la solution proposée}\label{sec:Limitations}
\subsection{Limitations de la méthodologie choisie}
\subsubsection{Différences entre les places minimums et les places prédites}
D'un point de vue holistique, l'hypothèse centrale de la méthodologie choisie est que le nombre de places de stationnement minimum requis est construit lors du développement de la propriété. Or dans certains cas, comme le logement unifamilial, cette hypothèse n'est pas nécessairement adéquate et d'autres méthodes (régression à partir de données réelles, apprentissage profond) seraient peut-être mieux à même de remplir ces enjeux pour certaines typologies de propriétés où ces hypothèses sont mal avisées. Les graphiques d'erreur présentés à la section
\subsection{Limitations de l'implémentation}
\subsubsection{Conversions d'unités}
\subsubsection{Capacité de faire des analyses alternatives}
\subsection{Limitations de données}
%%
%%  AMELIORATIONS FUTURES / FUTURE RESEARCH
%%
\section{Améliorations futures}
\subsection{Amélioration à la méthodologie par les règlements d'urbanisme}
\subsection{Validation statistique}
\subsection{Capacités à implémenter dans l'outil}
\subsection{Améliorations à l'outil}

